% Section III Bag of Tricks as in the shared Overleaf main.tex, 2026-09-14 (Shriram's version, mentor-owned).
\section{Bag of Tricks}
\label{sec:bag_of_tricks}

\subsection{Overview and Design Principle}

Vision--language--action (VLA) policies inherit the semantic generality of
large vision--language models, but also inherit their inference cost. This is
particularly restrictive in robotic manipulation, where inference is embedded
inside a feedback loop: an inefficient policy lowers the control frequency,
whereas an overly aggressive approximation can remove the visual or temporal
cues required for contact-rich behavior. We therefore seek training-free
interventions that reduce redundant computation while preserving the feedback
structure of the pretrained policy.

At time $t$, a VLA policy maps an RGB observation $I_t$ and language instruction
$\ell$ to an action or action chunk,
\begin{equation}
    A_t=\pi_\theta(I_t,\ell)
       =(a_{t,0},\ldots,a_{t,m-1}),
    \qquad a_{t,j}\in\mathbb{R}^{d_a}.
    \label{eq:vla_formulation}
\end{equation}
The parameters $\theta$ remain fixed throughout. Our central premise is that
VLA inference contains three distinct forms of redundancy: \emph{spatial
redundancy} in the visual observation, \emph{structural redundancy} across
decoder blocks, and \emph{temporal redundancy} across consecutive control
steps. Treating these as separate axes provides a principled way to determine
where computation may be removed and what failure mode each approximation can
introduce.

% This perspective also clarifies the scientific role of the individual methods.
Foveation and action repeat are intentionally simple controls that isolate
spatial and temporal perturbations. Depth pruning and VLA-Cache instantiate
established structural and cache-based acceleration principles
~\cite{shortgpt,efficientvla,vlacache}. Our main methodological contribution is
to connect these axes through two conservative mechanisms: (i) a guarded
action-reuse rule that interprets policy evaluation as an event-triggered
decision, and (ii) an interaction-aware temporal selector whose invariance mask
can be shared by representation fusion and model-internal caching.
% This unified
% view turns a collection of heuristics into a controlled study of when a VLA
% must perceive, reason, and act again.

\subsection{Foveation: Spatial Allocation of Perceptual Fidelity}
\label{sec:method_foveation}

Uniform image fidelity is expensive but not necessarily uniformly useful.
Manipulation tasks often concentrate geometric and semantic evidence around a
small set of objects, the gripper, and their interaction region. Motivated by
foveated perception~\cite{schwartz,traver}, we study whether a pretrained VLA
remains functional when visual fidelity decreases with distance from a central
region of interest.

For a keep ratio $\rho$, the sharp central disc has radius
$r_s=\sqrt{\rho HW/\pi}$. Let $d(x)$ denote the distance of pixel $x$ from the
image center and let $\tau(x)$ increase from zero at $r_s$ to one at the image
boundary. The foveated observation is a radial mixture
\begin{equation}
\begin{aligned}
\widetilde I_t(x)
&= w_0(x)I_t(x) + w_3(x)G_3(I_t)(x) \\
&\quad + w_9(x)G_9(I_t)(x), \\
&\text{s.t.}\quad
w_0(x)+w_3(x)+w_9(x)=1.
\end{aligned}
\label{eq:foveation}
\end{equation}
where $G_3$ and $G_9$ are Gaussian filters of increasing strength, and the
weights vary monotonically with $\tau(x)$. Pixels inside the fovea are retained
exactly. We evaluate keep ratios of $20\%$ and $50\%$.

% The purpose of this intervention is conceptual as much as computational.
% It

This tests the hypothesis that manipulation policies require high-frequency detail
only in a restricted spatial support. However, because the image resolution
and token count remain unchanged, fixed foveation does not itself reduce VLA
FLOPs. It is therefore a perceptual robustness baseline and a diagnostic for
future adaptive or task-conditioned foveation, rather than a claimed
acceleration method.

\subsection{Action Repeat: A Temporal Redundancy Control}
\label{sec:method_action_repeat}

Consecutive robot observations are strongly correlated, and the corresponding
actions often vary smoothly. The simplest exploitation of this redundancy is
to hold a predicted action for multiple control transitions. For repeat factor
$k$, the controller applies
\begin{equation}
    u_{t+j}=a_t,\qquad j=0,\ldots,k-1,
    \label{eq:action_repeat}
\end{equation}
and does not query the policy during the repeated steps. For an action-chunking
policy, the same operator may be applied to each element of the chunk.

Action repeat exposes a fundamental tension between inference efficiency and
feedback. It reduces the number of policy calls by approximately a factor of
$k$, but proportionally enlarges the open-loop interval. This is especially
problematic near contact, grasp closure, and final placement, where small
perceptual changes can require qualitatively different actions. We therefore
use $k\in\{2,4\}$ as negative temporal controls. Their role is to measure how
much apparent efficiency can be obtained by ignoring state-dependent risk and,
in turn establish a conditional reuse mechanism.

\subsection{Depth Pruning: Structural Redundancy in the Decoder}
\label{sec:method_depth_pruning}

The language-model decoder is a major source of per-query cost in OpenVLA.
Residual architectures suggest a natural redundancy criterion: if a block
produces only a small directional change in its hidden state over the target
distribution, then bypassing that block may preserve the local representation
while reducing both prefill and autoregressive computation.
% This reasoning
% follows Block Influence in ShortGPT and its adaptation to VLA acceleration
% ~\cite{shortgpt,efficientvla}.

For decoder block $j$, calibration example $n$, and token position $q$, we
measure
\begin{equation}
    s_j=\frac{1}{N}\sum_{n=1}^{N}
        \left[1-\frac{1}{Q_n}\sum_{q=1}^{Q_n}
        \cos\!\left(h^{\mathrm{in}}_{j,n,q},
                    h^{\mathrm{out}}_{j,n,q}\right)\right].
    \label{eq:block_influence}
\end{equation}
Low-influence blocks are candidates for removal. Selection is deliberately
conservative: the early decoder, which establishes multimodal processing, and
the final block, which directly precedes action generation, are protected;
adjacent blocks are not removed together; and the selected indices are frozen
before evaluation. The resulting network is structurally shortened rather
than emulated with identity layers, so the method removes real computation and
maintains a consistent generation cache.

The underlying assumption is distributional rather than universal. A block
that is nearly redundant on calibration trajectories may become important
under a different scene, instruction, or contact state. Moreover, cosine
similarity measures representational direction, not action equivalence.
% Accordingly, we evaluate increasing removal budgets and separately inspect
% continuous-action and binary-gripper disagreement.
The contribution here is
not a new pruning score, but a VLA-specific, cache-safe, calibration-controlled
application of structural redundancy.

\subsection{Guarded Reuse: State-Conditional Policy Evaluation}
\label{sec:method_guarded_reuse}

Unconditional repetition treats every control phase as equally predictable.
Our guarded-reuse mechanism instead asks whether the information available at
the current step is sufficiently consistent with continued execution of the
previous action. This can be interpreted as event-triggered inference: the
expensive VLA is evaluated only when a perceptual or control event invalidates
the local assumption of temporal continuity.

The gate combines complementary evidence. A global image-change statistic
detects camera or scene motion; a maximum local-patch statistic detects small
changes that would disappear in a whole-image average; agreement between the
two most recent inferred pose actions estimates local policy smoothness; and a
gripper-consistency test protects discrete mode transitions. We additionally
reject reuse for near-zero translational commands.
% Although such commands look
% stable numerically, they frequently occur during precise alignment or contact,
% where fresh feedback is most valuable.

Let $d_t^g$ and $d_t^l$ be global and maximum-local differences between cheap
signatures of consecutive observations, let $c_t$ be the cosine similarity of
the two most recent densely inferred six-dimensional pose actions, let $v_t$
be the current translation magnitude, and let $r_t$ count consecutive reuses.
The decision is
\begin{equation}
\begin{aligned}
G_t ={}&
[d_t^g \le \tau_g] \wedge [d_t^l \le \tau_l]
\wedge [c_t \ge \tau_a] \wedge [v_t \ge \tau_v] \\
&\wedge [\text{gripper state unchanged}]
\wedge [r_t < R_{\max}], \\[2pt]
u_t ={}&
\begin{cases}
u_{t-1}, & G_t = 1,\\
\pi_\theta(I_t,\ell), & G_t = 0.
\end{cases}
\end{aligned}
\label{eq:guarded_reuse}
\end{equation}
The current observation is examined at every step even when inference is
skipped. Any failed condition immediately returns the controller to dense
inference, and a maximum reuse horizon prevents prolonged open-loop behavior.

We study strict, moderate, and aggressive operating points that jointly relax visual-change and action-agreement thresholds, tracing the efficiency–reliability trade-off. Unlike action-reuse methods such as FlashVLA~\cite{flashvla}, our conservative guard combines global and local visual evidence, action geometry, and gripper state to determine when action similarity is sufficiently reliable to skip a policy query.

\subsection{VLA-Cache: Task-Aware Reuse of Visual Computation}
\label{sec:method_vla_cache}

Guarded reuse skips the entire model, whereas VLA-Cache~\cite{vlacache} targets
redundancy inside a policy call. Its premise is that consecutive manipulation
frames contain many visually stable patches whose previously computed states
can be retained. Crucially, visual stability is not equivalent to task
irrelevance: a target object may remain static while still determining the
correct action. VLA-Cache therefore combines inter-frame similarity with
language-conditioned attention.

For patch $i$, let $\gamma_{t,i}$ denote its cosine similarity across frames
and let $r_{t,i}$ be its text-to-vision attention score from the preceding
query. Reuse is restricted to patches that are both sufficiently stable and
outside the most task-relevant set:
\begin{equation}
    \mathcal C_t=
    \operatorname{TopK}\!\left({i:\gamma_{t,i}\ge\tau_s\},K_s\right)
    \setminus\operatorname{TopK}(r_t,K_r).
    \label{eq:vla_cache}
\end{equation}
The corresponding historical cache entries are reused according to a
layer-adaptive schedule derived from attention concentration. This preserves
historical information rather than simply discarding tokens, distinguishing
caching from token pruning.

% We include VLA-Cache as a recent literature baseline because it directly tests
% whether temporal visual redundancy produces actual decoder-level savings.
The first observation remains dense, while stability and relevance filtering may retain fewer patches than the nominal budget. We therefore separate cache construction from steady-state cost and report actual reuse. This motivates whether task-conditioned temporal invariance can govern both representation fusion and computational reuse.


\subsection{Interaction-Aware Temporal Fusion}
\label{sec:method_temporal_fusion}

Frame-independent VLA inference discards a useful prior: physical scenes evolve
continuously, so a stable visual representation should not fluctuate
arbitrarily between adjacent steps. Temporal fusion exploits this continuity by
retaining historical projected tokens for regions judged invariant. Our formulation develops this reasoning with
an explicit interaction-aware protection rule and a cache-compatible shared
mask.

For every visual patch $i$, we compute three signals. Motion $m_{t,i}$ captures
physical change; normalized image entropy $e_{t,i}$ acts as a task-agnostic
proxy for local structure and object boundaries; and optional relevance
$r_{t,i}$ captures language-conditioned importance from the preceding query.
These signals are complementary rather than interchangeable. Motion protects
dynamic regions, entropy protects visually informative regions even before
their role is known, and attention protects instruction-relevant regions that
may be static.
% Since patch boundaries and attention maps are imperfect, the
% protected mask is spatially dilated to preserve the neighborhood of potential
% interactions.
Let $\mathcal P_t$ be the union of high-motion, high-entropy, task-relevant,
and externally protected patches after dilation. Among its complement, we rank
patches by the combined risk $R_{t,i}=m_{t,i}+e_{t,i}+r_{t,i}$ and select a
bounded reusable set
\begin{equation}
\begin{aligned}
\mathcal P_t
&=\operatorname{Dilate}\!\left(
    \{i:m_{t,i}>\tau_m\}
    \cup \operatorname{TopFrac}(e_t,\eta_e)\right.\\
&\qquad\left.
    \cup \operatorname{TopFrac}(r_t,\eta_r)
    \cup \mathcal P_t^{\mathrm{interaction}}
    \right),\\
\mathcal R_t
&=\operatorname{BottomK}\!\left(
    \{R_{t,i}:i\notin\mathcal P_t\},
    \lfloor\beta P\rfloor
    \right).
\end{aligned}
\label{eq:interaction_selector}
\end{equation}
This construction encodes a conservative asymmetry: evidence is required to
reuse a patch, whereas any plausible indication of motion, information content,
or task relevance is sufficient to recompute it.

Let $z_{t,i}$ be the current projected token and
$\widetilde z_{t-1,i}$ the previously retained representation. Temporal fusion
updates the decoder input as
\begin{equation}
    \widetilde z_{t,i}=
    \begin{cases}
      \widetilde z_{t-1,i},&i\in\mathcal R_t,\\
      z_{t,i},&i\notin\mathcal R_t.
    \end{cases}
    \label{eq:temporal_fusion}
\end{equation}
% Because historical fused tokens are propagated recursively, reuse can suppress
% transient visual noise but can also accumulate stale information. Periodic
% dense keyframes bound this temporal horizon, while an abrupt-motion detector
% can force an early reset. We evaluate motion--entropy, task-aware, and
% conservative-adaptive variants to isolate the contribution of semantic
% relevance and reuse aggressiveness.

% Temporal fusion alone does not reduce model FLOPs: the current visual encoder,
% projector, and decoder are still executed, and the token sequence is unchanged.
% Its role is to test whether temporally coherent representations improve policy
% stability. The acceleration opportunity arises from sharing $\mathcal R_t$
% with the cache path, so that a patch considered invariant at the representation
% level is also eligible for internal compute reuse. This \emph{shared invariance
% mask} is the central conceptual link between temporal fusion and caching. It
% prevents two independently tuned mechanisms from making contradictory decisions
% about interaction-critical regions.
Recursive fusion can suppress transient noise but risks stale information; periodic keyframes and abrupt-motion resets bound this effect. We evaluate motion--entropy, task-aware, and conservative-adaptive variants to separate semantic relevance from reuse aggressiveness.

Temporal fusion itself does not reduce FLOPs, since the full policy still executes. Efficiency arises by sharing $\mathcal R_t$ with the cache, making representation-level invariance directly determine compute reuse. This \emph{shared invariance mask} unifies fusion and caching while protecting interaction-critical regions.

% The current implementation exposes this shared mask, but selective KV/KQV
% execution under the optimized SDPA path remains future integration work.
% Accordingly, we evaluate temporal fusion and VLA-Cache separately and do not
% attribute a measured acceleration to the combined interface.
